\section*{Chapitre \ref{ch:nw:programmable-hardware-ii-trading-designs} --- Matériel programmable II~: architectures de trading}

\begin{solution}{exo:nw:programmable-hardware-ii-trading-designs:1}
Le prix se termine à l'octet $20 + 2 + 35 = 57$, dans le mot 7 (octets 56 à 63). L'ordre est enregistré deux cycles plus tard, à la fin du cycle 9~:
dix cycles depuis le début du premier mot du paquet, $10 \times 6.4 = \qty{64}{\nano\second}$.
\end{solution}

\begin{solution}{exo:nw:programmable-hardware-ii-trading-designs:2}
Un jeton toutes les 64 cycles~: $156.25 \times 10^6 / 64 \approx 2.44$ millions d'ordres par seconde en régime établi, après une rafale initiale de 4.
\end{solution}

\begin{solution}{exo:nw:programmable-hardware-ii-trading-designs:3}
$50 \times 156.25 = 7\,812.5$ cycles~; $4 + \lfloor 7\,812.5 / 64\rfloor = 4 + 122 = 126$ ordres.
\end{solution}

\begin{solution}{exo:nw:programmable-hardware-ii-trading-designs:4}
Les 47 messages qualifiants supplémentaires (160 contre 113) arrivent par rafales, lorsque les cours vendeurs proches du seuil sont émis~: le seau à jetons, déjà
vidé par les premiers ordres de chaque rafale, en rejette 22 et en laisse passer 25.
\end{solution}

\begin{solution}{exo:nw:programmable-hardware-ii-trading-designs:5}
Un ordre ramené à la limite n'est pas l'ordre que la stratégie a calculé~; l'écrêtage modifie le risque en silence et masque une erreur de configuration ou de modèle.
Le rejet garde le comportement du matériel simple et auditable, et rend l'erreur visible dans le compteur de rejets.
\end{solution}

\begin{solution}{exo:nw:programmable-hardware-ii-trading-designs:6}
Le seuil lui-même (un modèle qui utilise de nombreuses entrées et l'historique), la position et l'exposition sur l'ensemble des plateformes (un état partagé avec d'autres systèmes),
et la décision même de trader (actualité, suspensions de cotation, limites du desk). Chacun demande un état, une arithmétique ou un jugement, ce qui coûterait de la logique et des semaines de
vérification pour chaque modification.
\end{solution}

\begin{solution}{exo:nw:programmable-hardware-ii-trading-designs:7}
Un seul ordre~: le montage d'essai dure moins de 10\,000 cycles, si bien que le seau à jetons ne se recharge jamais après son unique jeton. Avec le bit kill positionné un cycle
avant le dixième ordre de l'exécution sans restriction, exactement neuf ordres sortent.
\end{solution}

\begin{solution}{exo:nw:programmable-hardware-ii-trading-designs:8}
Les 13 nanosecondes (chiffre de référence) courent du dernier bit nécessaire à la décision jusqu'au premier bit de l'ordre. La latence wire-to-wire ajoute le
temps d'arrivée des bits déclencheurs, les transceivers de la carte, et la fibre et les équipements de la cage dans les deux sens~: des centaines de nanosecondes dans
la cage du chapitre.
\end{solution}

\begin{solution}{pb:nw:programmable-hardware-ii-trading-designs:1}
\begin{enumerate}
\item $263 + 120 + 19.2 + 120 + 295 \approx \qty{817}{\nano\second}$~: \qty{0.82}{\micro\second}.
\item Trois cycles, \qty{19.2}{\nano\second}~: environ 2~\%.
\item \qty{0.82}{\micro\second}~: chaque étape du modèle est une constante (fibre, latence fixe des transceivers, cycles fixes).
\item D'ici \qty{3.57}{\micro\second} à la médiane et \qty{34.5}{\micro\second} au 99\up{e}~centile.
\item 75~\% pour le chemin matériel, 0.01~\% pour le chemin logiciel.
\item 100~\% et 57~\%.
\item La fibre de la cage et les transceivers~: ce qui reste relève de la physique (distance) et de la sérialisation, des chapitres~9 à~14.
\item La corrélation entre le concurrent et l'entreprise (les deux voient les mêmes rafales), la mise en file d'attente à la passerelle de la plateforme, et tout sauf une course isolée.
\item $4 + \lfloor 3\,125/64\rfloor = 52$ ordres.
\item $4 + \lfloor 16/64\rfloor = 4$ ordres~: la rafale du seau.
\item Pour qu'un ordre déjà passé l'étape de contrôle des risques soit aussi arrêté dans son dernier cycle~: le bit kill bloque alors tout à partir du cycle où il
  est positionné.
\item Sur la carte~: taille, tunnel de prix par rapport à une référence, débits d'ordres et de notionnel, le bit kill. Dans la couche du courtier, sous son contrôle~:
  seuils de crédit et de capital, et son propre kill switch.
\item \textbf{Résultat clé.} \qty{0.82}{\micro\second} en wire-to-wire dans le matériel, contre une médiane de \qty{4.39}{\micro\second} dans le logiciel~; face à un
  concurrent à \qty{1}{\micro\second}, ils gagnent 75~\% et 0.01~\% des courses~; un kill activé au bout de \qty{20}{\micro\second} laisse sortir au plus 52 ordres.
\item Guère~: à \qty{5}{\micro\second}, le chemin logiciel gagne déjà 57~\% et le reste se répartit dans la queue de la distribution~; mieux vaut corriger cette queue.
\item Flexibilité, rapidité de modification et l'essentiel de l'expressivité de la stratégie~; elle gagne en revanche une charge de vérification.
\item Aussi souvent que le logiciel peut le calculer et l'écrire, en quelques microsecondes~; le logiciel de la stratégie en est responsable, dans les limites de la couche de risque.
\item Pour connaître sa position et la rapprocher des accusés de réception de la plateforme~: la carte agit seule, et l'entreprise ne doit pas perdre la trace de
  ce qu'elle a fait.
\item Rejouer le flux enregistré dans le HDL en simulation et sur la carte dans un \omterm{def:nw:programmable-hardware-i-architecture-and-toolchain:tb}{banc de test}, comparer les deux aux modèles, et lancer
  la certification de la plateforme.
\item Seuils ou limites périmés~: la conception doit conserver un numéro de version et un heartbeat issus du logiciel, et cesser de déclencher lorsqu'ils sont
  périmés.
\item Le matériel décide du quand~; le logiciel décide du quoi, et même s'il faut agir.
\end{enumerate}
\end{solution}

\begin{solution}{iq:nw:programmable-hardware-ii-trading-designs:1}
Il traite un nombre fixe d'octets par cycle, avec une machine à états qui suit la position dans l'en-tête et dans chaque message à préfixe de
longueur, en capturant les champs au fur et à mesure que passent leurs octets~; les champs d'un message sont connus au cycle où arrive son dernier octet, et les décisions peuvent commencer
avant que le paquet ne se termine ou ne soit vérifié.
\iqlookfor{état en flux, sans mise en mémoire tampon, et décision avant la fin de la trame.}
\end{solution}

\begin{solution}{iq:nw:programmable-hardware-ii-trading-designs:2}
Taille, tunnel de prix, limites de débit (un seau à jetons), détection des doublons et bit kill, chacun étant une comparaison ou un compteur placé dans une étape du pipeline que
tout ordre doit traverser. On le prouve par construction (le seul chemin vers la sortie passe par l'étape), par simulation de chaque frontière et
en gardant les contrôles du courtier indépendants de la configuration de l'entreprise.
\iqlookfor{un chemin unique vers le réseau et des contrôles indépendants.}
\end{solution}

\begin{solution}{iq:nw:programmable-hardware-ii-trading-designs:3}
Modèles, paramètres, positions, risque inter-plateformes et tout ce qui exige un historique ou un jugement~: le logiciel calcule ce qui vaudrait la peine d'être fait,
le matériel l'applique au moment où le marché l'offre. Les modifications matérielles sont lentes et coûteuses à vérifier~; celles du logiciel sont peu coûteuses.
\iqlookfor{répartition entre quoi et quand, et coût du changement.}
\end{solution}

\begin{solution}{iq:nw:programmable-hardware-ii-trading-designs:4}
Avec un \omterm{def:nw:programmable-hardware-ii-trading-designs:toe}{moteur de déchargement TCP} (TOE) sur la carte, qui gère les numéros de séquence de la connexion, les accusés de réception et les retransmissions, et un \omterm{def:nw:programmable-hardware-ii-trading-designs:template}{gabarit d'ordre}
dont les en-têtes TCP et IP et les sommes de contrôle sont mis à jour à l'émission~; la couche de session (logon, heartbeats) est partagée avec le logiciel.
\iqlookfor{état de connexion dans le matériel et un gabarit à champs mis à jour.}
\end{solution}

\begin{solution}{iq:nw:programmable-hardware-ii-trading-designs:5}
Exécuter les deux sur les mêmes entrées enregistrées et synthétiques, dans deux simulateurs, et comparer chaque sortie à son cycle~; ajouter des contrôles de référence
au niveau des messages (chaque message qualifiant est décidé une seule fois). Un désaccord est un bogue dans l'un des deux jusqu'à preuve du contraire~: réduire l'entrée au
plus petit cas possible et lire le chronogramme.
\iqlookfor{comparaison au niveau du cycle et réduction méthodique.}
\end{solution}

\begin{solution}{iq:nw:programmable-hardware-ii-trading-designs:6}
Depuis quel bit jusqu'à quel bit elle est mesurée, pour quelles tailles de trame et quels types de messages, si elle inclut les transceivers et le
TCP de l'ordre, sous quelle charge et avec quels contrôles de risque activés, et comment le chiffre se rapporte à mon chemin wire-to-wire.
\iqlookfor{la définition de la mesure.}
\end{solution}
